% Introduction (Kurose 5-paragraph structure).

\section{Introduction}
\label{sec:intro}

Low Earth orbit (LEO) mega-constellations have become a mainstream component of
the global Internet: hundreds to thousands of satellites at a few hundred
kilometers of altitude now deliver low-latency broadband to regions that
terrestrial infrastructure cannot economically reach~\cite{smallsats2020,
handley:ground-relays:2019}.
Their low altitude yields round-trip times competitive with, and sometimes
better than, long-haul fiber, which has driven intense study of constellation
design and performance at scale~\cite{kassing2020exploring,lai2020starperf}.
As these systems grow to thousands of nodes carrying production traffic, the
network layer faces a defining challenge: it must keep packets flowing while the
physical topology is in perpetual, rapid motion.

Because satellites orbit continuously, a LEO network is most naturally described
as a \emph{time-varying graph} in which inter-satellite and ground--satellite
links (GSLs) appear and disappear on a schedule dictated by orbital
mechanics~\cite{seong:topological-design:1995,kassing2020exploring}.
A ground station must therefore hand its uplink from a setting satellite to a
rising one at recurring, predictable instants, a GSL \emph{handover}.
Distributed link-state routing such as open shortest path first
(OSPF)~\cite{sidhu:ospf:1993}, the de facto baseline in container-based LEO
emulators~\cite{lai2023starrynet}, learns each mutation reactively through
hello timeouts and link-state flooding, and reconverges only \emph{after} the
link change has completed.
During that interval, traffic is forwarded along stale or missing next-hops and
is silently black-holed.
This raises the central question of this paper: since handover times are known
in advance from orbital prediction, can routing exploit that knowledge to keep
the data plane continuously connected across topology changes?

In this paper we present OrbitGraph, a graph-centric, software-defined routing
scheme for LEO constellations.
OrbitGraph models each orbital snapshot as a delay-weighted graph $G_t$, computes
all-pairs shortest paths centrally with Dijkstra's algorithm, and programs the
kernel forwarding information base (FIB) of each node by pushing only the routes
that changed since the previous snapshot.
At a handover, the controller installs the post-handover FIB \emph{before} the
old GSL is torn down, turning orbital predictability into make-before-break
forwarding at the topology level.
We implement OrbitGraph~\cite{online:orbitgraph:2025} on the StarryNet emulator~\cite{lai2023starrynet}, which
runs real TCP/IP stacks and an OSPF control plane in containers, and evaluate it
against that distributed baseline on Walker constellations from 27 to 102 nodes
using a continuous outage probe aligned to routing events.
Our contributions are fourfold: a graph-centric software-defined networking
(SDN) control model that unifies initialization, delay refresh, link failure,
and GSL handover as shortest-path computations over a sequence of snapshots
$G_t$; an incremental, make-before-break FIB install that programs only changed
next-hops, reducing handover updates by roughly two orders of magnitude relative
to a full-table push; a proactive handover mechanism that commits the new
forwarding state while both contacts briefly coexist; and an empirical study
showing 0\,ms data-plane outage on coverage-preserving GSL handovers where OSPF
incurs multi-second black holes, alongside millisecond-range shortest-path
recomputation, together with a characterization of the boundaries of the
approach: a satellite coverage gap that severs the path for both protocols, and
unplanned bulk failure where centralized install does not yet outperform
distributed reconvergence.

OrbitGraph differs from prior LEO routing along a single axis: \emph{what}
forwarding state to install and \emph{when} relative to link lifetime.
Snapshot and offline schemes precompute per-state tables and load them on
board~\cite{seong:topological-design:1995}; distributed protocols optimize hop
count or convergence speed without a global view~\cite{gregory:satellite-routing:2022,
pam:opspf:2019}; and satellite SDN research has largely targeted controller
\emph{placement} rather than route computation~\cite{papa:sdn-cpp:2018,
guo:spda:2022}.
To our knowledge, none of these install real kernel forwarding state on emulated
satellite nodes and measure end-to-end availability through handover under
proactive update ordering; Section~\ref{sec:related} compares these lines in
detail.

The remainder of this paper reviews software-defined and graph-centric
networking principles, situates OrbitGraph among prior LEO routing work,
describes the StarryNet emulation platform and OrbitGraph design, reports the
scaling methodology and results, discusses limitations, and concludes.
